% =====================================================================
%  PLANTILLA DE PROBLEMA
%  1. Copia este archivo:  problemas/pN.tex
%  2. Agrégalo con \input{problemas/pN} en el archivo
%     de su parte (partes/parteX_...tex).
%  3. (Opcional) Crea python/problemas/pN.py para sus
%     figuras y tablas; se guardan en figuras/ y tablas/ con prefijo pN_.
% =====================================================================
\problema{NN}{Título del problema}{0}

\begin{enunciado}
Texto del enunciado.
\begin{enumerate}[label=(\alph*)]
  \item Primer inciso.
  \item Segundo inciso.
\end{enumerate}
\end{enunciado}

\paragraph{Referencias en los apuntes.} Sesión XX (...).

\inciso{a}{Título corto del inciso}

\paso{1}{Generación del mallado}
Desarrollo.

\begin{porque}
Explicación de por qué se hace este paso.
\end{porque}

\resultado{u_{\ell+1} = \dots}

\begin{ojo}
Error común o detalle que confunde.
\end{ojo}

\inciso{b}{Título corto del inciso}
Desarrollo.

% Material generado por Python:
% \tablapython{pN_tabla}
% \figurapython{pN_figura}{Pie de figura.}{fig:pN}{0.75}
